\section*{Chapitre \ref{ch:g7:central} --- Symétrie centrale et parallélogrammes}

\begin{solution}{exo:g7:central:1}
Image par demi-tour~: $M'$ est à quatre carreaux à \emph{gauche} de
$O$ et un carreau en \emph{bas} (les deux directions inversées). Image
par réflexion par rapport à la droite verticale~: quatre carreaux à
gauche mais toujours un carreau en \emph{haut}. Les deux images
diffèrent~: elles sont des images miroir verticales l'une de l'autre.
\end{solution}

\begin{solution}{exo:g7:central:2}
$A'B' = AB$ (les demi-tours conservent les longueurs) et
$(A'B') \parallel (AB)$ (l'image d'une droite est une droite
\omterm{def:g6:lines:perp}{parallèle}, \cref{prop:g7:central:preserved}).
\end{solution}

\begin{solution}{exo:g7:central:3}
Chiffres de calculatrice avec un \omterm{def:g7:central:center}{centre de symétrie}~: $0$, $2$, $5$,
$8$ (et $1$ s'il est dessiné comme une simple barre). Lettres~: H, N,
S, Z, O en ont un~; A et T n'en ont pas (elles ont un axe à la
place).
\end{solution}

\begin{solution}{exo:g7:central:4}
Côtés opposés égaux~: $MN = KL = 7$ cm et $NK = LM = 4$ cm. Angles
opposés égaux~: $\widehat M = \widehat K = 115^\circ$. Les diagonales
se coupent en leur milieu commun, et ce point est $O$~: le milieu de
$[LN]$ est $O$.
\end{solution}

\begin{solution}{exo:g7:central:5}
Tracer un \omterm{def:g6:lines:objects}{segment} $[AC]$ de $8$ cm, marquer son milieu $O$, tracer
par $O$ une droite faisant $60^\circ$ avec $(AC)$, et placer $B$ et
$D$ dessus à $2.5$ cm de $O$ de part et d'autre. Joindre $A$, $B$,
$C$, $D$~: les diagonales se coupent en leurs milieux, donc $ABCD$
est un \omterm{def:g7:central:parallelogram}{parallélogramme} par définition.
\end{solution}

\begin{solution}{exo:g7:central:6}
Les angles consécutifs sont supplémentaires~: à côté de $115^\circ$
se trouve $180 - 115 = 65^\circ$. Les angles opposés sont égaux, donc
les quatre angles sont $115^\circ$, $65^\circ$, $115^\circ$,
$65^\circ$ (\omterm{def:g1:addition:def}{somme} $360^\circ$).
\end{solution}

\begin{solution}{exo:g7:central:7}
De $A(1,1)$ à $O(2.5, 2.5)$~: $1.5$ à droite et $1.5$ en haut~; en
continuant de même~: $A'(4, 4)$. De $B(4,2)$ à $O$~: $1.5$ à gauche
et $0.5$ en haut~; en continuant~: $B'(1, 3)$.
\end{solution}

\begin{solution}{exo:g7:central:8}
Non. Un trapèze \omterm{def:g6:shapes:triangles}{isocèle} a ses deux côtés obliques égaux sans être un
\omterm{def:g7:central:parallelogram}{parallélogramme} (les deux côtés \omterm{def:g6:lines:perp}{parallèles} ont des longueurs
différentes). Ce qui suffit est le critère 3 de
\cref{thm:g7:central:recognize}~: deux côtés opposés égaux
\emph{et \omterm{def:g6:lines:perp}{parallèles}}.
\end{solution}

\begin{solution}{exo:g7:central:9}
$O$ est le milieu de $[BC]$ par choix, et il est le milieu de
$[AA']$ par construction du \omterm{def:g7:central:def}{symétrique}. Le \omterm{def:g4:geometry:polygon}{quadrilatère} $ABA'C$ a
des diagonales $[AA']$ et $[BC]$ partageant le milieu $O$~: le
critère 1 (la définition) en fait un \omterm{def:g7:central:parallelogram}{parallélogramme} sans coût
supplémentaire.
\end{solution}

\begin{solution}{exo:g7:central:10}
«~Angles opposés égaux $\Rightarrow$ \omterm{def:g7:central:parallelogram}{parallélogramme}~»~: \emph{vrai}
(c'est un autre critère de reconnaissance~; avec la \omterm{def:g1:addition:def}{somme} des angles
$360^\circ$, des paires opposées égales forcent des angles
consécutifs supplémentaires, d'où des côtés \omterm{def:g6:lines:perp}{parallèles} par
\cref{thm:g7:triangles:equalities}).

«~Diagonales \omterm{def:g6:lines:perp}{perpendiculaires} $\Rightarrow$ \omterm{def:g6:shapes:quadrilaterals}{losange}~»~: \emph{vrai}
pour un \omterm{def:g7:central:parallelogram}{parallélogramme} (les diagonales se coupent déjà en leurs
milieux).

«~Un \omterm{def:g3:shapes:rightangle}{angle droit} $\Rightarrow$ rectangle~»~: \emph{vrai} --- les
angles consécutifs sont supplémentaires, donc les quatre deviennent
droits.
\end{solution}

\begin{solution}{exo:g7:central:11}
Quel que soit le \omterm{def:g4:geometry:polygon}{quadrilatère} --- même très irrégulier --- les
milieux $I$, $J$, $K$, $L$ forment toujours un
\emph{\omterm{def:g7:central:parallelogram}{parallélogramme}}. La preuve avec le théorème des milieux est au
\cref{ch:g8:midpoints} (\cref{exo:g8:midpoints:9}).
\end{solution}

\begin{solution}{pb:g7:central:1}
\textbf{1.} $M(3, 1) \to (-3, -1)$~; $N(-2, 4) \to (2, -4)$~;
$P(0, -5) \to (0, 5)$. Règle~: le demi-tour de centre l'origine
envoie $(x, y)$ sur $(-x, -y)$ --- les deux coordonnées changent de
signe.

\textbf{2.} À partir de $M(5, 3)$~: trois carreaux à droite de $C$
deviennent trois carreaux à gauche, deux vers le haut deviennent
deux vers le bas~: image $(-1, -1)$. Vérification de la recette~: le
double du centre moins le point,
$(2 \times 2 - 5,\ 2 \times 1 - 3) = (-1, -1)$.

\textbf{3.} Le demi-tour de centre la carte échange les points deux
par deux. Avec un \omterm{def:g3:numbers:evenodd}{nombre impair} de points, l'un d'eux n'a pas de
partenaire~: il est donc sa propre image, et le seul point qui soit
sa propre image est le centre lui-même. D'où le point du milieu du
3, du 5, du 7\,\dots\ pile au centre (regarder une vraie carte~: il
y est).

\textbf{4.} Par $O_1(2,0)$~: $A(1,1) \to (3,-1)$~; puis par
$O_2(5,0)$~: $(3,-1) \to (7,1)$. Départ $A(1,1)$, arrivée
$(7,1)$~: glissé de $6$ vers la droite, sans retournement. De même
pour $B(0,3) \to (4,-3) \to (6,3)$~: glissé de $6$ vers la droite.
Les deux demi-tours se ramènent à une \emph{translation}.

\textbf{5.} Le glissement vaut $6$, et $O_1 O_2 = 3$~: la
translation parcourt le \emph{double} de la distance entre les
centres (dans le sens de $O_1$ vers $O_2$) --- exactement comme les
deux miroirs \omterm{def:g6:lines:perp}{parallèles} translataient du double de leur écart dans
\cref{pb:g6:symmetry:1}.

\textbf{6.} Le demi-tour de centre $O$ envoie le \omterm{def:g7:central:parallelogram}{parallélogramme} sur
lui-même (c'est ce que signifie pour lui «~\omterm{def:g7:central:center}{centre de symétrie}~»,
\cref{thm:g7:central:properties}). Une droite passant par $O$ est
elle aussi envoyée sur elle-même~: les deux points où elle coupe le
bord s'échangent donc, et ils sont \omterm{def:g7:central:def}{symétriques} par rapport à $O$.
Les deux morceaux du \omterm{def:g7:central:parallelogram}{parallélogramme} de part et d'autre de la droite
s'échangent également~: ce sont donc des copies exactes
(\cref{prop:g7:central:preserved}) --- de même \omterm{def:g6:measure:area}{aire}.

\textbf{7.} Couper suivant la droite qui joint les \emph{deux
centres}~: celui du rectangle du gâteau et celui du trou. Par la
question 6 appliquée au rectangle du gâteau, la coupe partage le
gâteau en deux~; appliquée au rectangle du trou (la droite passe
aussi par son centre), elle partage le trou en deux. Chaque part
reçoit la \omterm{def:g3:division:half}{moitié} du gâteau moins la \omterm{def:g3:division:half}{moitié} du trou~: parfaitement
équitable.

\textbf{8.} $D = (0 + 8 - 6,\ 0 + 5 - 1) = (2, 4)$ (les diagonales
ont le même milieu, donc $D$ est l'image de $B$ par le demi-tour de
centre le centre du \omterm{def:g7:central:parallelogram}{parallélogramme}). Centre~: milieu de $[AC]$,
$O(4,\ 2.5)$.

\textbf{9.} Dans le \omterm{def:g7:central:parallelogram}{parallélogramme} $ABA'C$, les côtés opposés sont
égaux~: $BA' = AC$. L'inégalité triangulaire dans $ABA'$ donne
$AA' < AB + BA' = AB + AC$. Comme $O$ est le milieu de $[AA']$, on a
$AA' = 2\,AO$, donc $2\,AO < AB + AC$, c'est-à-dire
$AO < \frac{AB + AC}{2}$.

\textbf{10.} Majoration~: $AO < \frac{5 + 7}{2} = 6$ cm.
Minoration~: dans $ABA'$, $AA' > BA' - AB = 7 - 5 = 2$, donc
$AO > 1$ cm. La médiane issue de $A$ est strictement comprise entre
$1$ et $6$ cm.

\textbf{11.} \omterm{def:g6:lines:objects}{Segment}~: oui, son milieu. Droite~: oui --- chacun de
ses points est un centre. \omterm{def:g6:shapes:triangles}{Triangle équilatéral}~: aucun. Carré~: oui,
le point d'intersection des diagonales. Cercle~: oui, son centre.
S et Z~: oui, leur point milieu (retourner la page~: elles se lisent
de la même façon).

\textbf{12.} Le demi-tour apparierait les trois \omterm{def:g5:solids:def}{sommets} entre eux~;
trois étant \omterm{def:g3:numbers:evenodd}{impair}, un \omterm{def:g5:solids:def}{sommet} $A$ serait sa propre image, ce qui
obligerait $A$ à être le centre du demi-tour. Les deux autres
\omterm{def:g5:solids:def}{sommets} $B$ et $C$ seraient alors \omterm{def:g7:central:def}{symétriques} par rapport à $A$ ---
donc $A$ serait le milieu de $[BC]$, ce qui alignerait les trois
\omterm{def:g5:solids:def}{sommets}. Or un triangle n'a pas trois \omterm{def:g5:solids:def}{sommets} alignés~:
contradiction. Aucun centre n'existe.

\textbf{13.} Le même argument de parité~: un demi-tour conservant le
\omterm{def:g4:geometry:polygon}{polygone} apparie ses \omterm{def:g5:solids:def}{sommets}~; un \omterm{def:g3:numbers:evenodd}{nombre impair} de \omterm{def:g5:solids:def}{sommets} impose un
\omterm{def:g5:solids:def}{sommet} fixe, qui doit être le centre et le milieu du \omterm{def:g6:lines:objects}{segment} joignant
deux autres \omterm{def:g5:solids:def}{sommets} --- trois \omterm{def:g5:solids:def}{sommets} alignés, impossible dans un
\omterm{def:g4:geometry:polygon}{polygone}\,\dots\ donc aucun \omterm{def:g4:geometry:polygon}{polygone} ayant un \omterm{def:g3:numbers:evenodd}{nombre impair} de
\omterm{def:g5:solids:def}{sommets} n'a de \omterm{def:g7:central:center}{centre de symétrie}.

\textbf{14.} Effectuer le demi-tour de centre $O_1$, puis celui de
centre $O_2$, envoie la figure sur elle-même à chaque fois --- donc
leur composée aussi. D'après la question 5, cette composée est la
translation du double de la distance $O_1 O_2$, qui n'est pas nulle
puisque les centres diffèrent. Or une figure tenant sur une feuille
ne peut pas être égale à elle-même glissée d'une quantité fixe non
nulle (la glisser assez de fois la fait sortir entièrement de la
feuille). Une figure bornée a donc au plus un \omterm{def:g7:central:center}{centre de symétrie}.

\textbf{15.} Sur la bande sans fin, le demi-tour de centre le point
situé à mi-chemin entre une empreinte gauche et l'empreinte droite
suivante envoie le motif sur lui-même~; le point situé un pas
complet plus loin aussi~: voilà deux centres distincts. La composée
des deux demi-tours est la translation du double du demi-pas ---
exactement le glissement d'un pas qui, visiblement, envoie la frise
infinie sur elle-même, comme l'annonce la question 14. Les motifs non
bornés vivent selon d'autres règles~: c'est la mathématique des
papiers peints.
\end{solution}
